\begin{abstract}
We study the memory required to attain a statistical saddle in a causal
experiment. For finite interfaces, admissible registers give normalized
positive realizations of the experiment's continuation channels. An
unrestricted minimax saddle determines an exposed Bayes face; optimal
finite-memory realizations must also satisfy contact and stationarity
at the supporting parameters. These additional conditions can fix
responses on Bayes ties and thereby force memory that a Bayes best
response alone would not require. For a continuous family of marked
product experiments we obtain exactly five states at the decision cut
of an optimal two-preparation audit. For serial validation with the
candidate read first, the minimum whole-schedule peak is exactly twelve.
Both converses include stochastic encoders and continuations. They are
proved by positive face packing, not by minimizing a chosen deterministic
machine. We also identify the all-preparation limiting value and derive
the closest-product geometry directly from the experimental probabilities.
Every least-favorable prior localizes near two intrinsic contact points;
symmetric priors converge to their equal mixture with an explicit
quadratic-transport bound. Prepared predictive quotients, causal
simulation, informative finite-memory control, exact one- and
two-preparation saddles, and positive-noise physical realization are
developed within the same resource framework.
\end{abstract}
\maketitle
\tableofcontents

\section{Introduction}
A statistical saddle determines the best expected decision, but need not
determine a machine that can make that decision. A Bayes response can
lose information that is immaterial under its prior yet necessary for
uniform optimality over the parameter family. In a causal experiment,
this distinction has a resource cost: discarded information cannot be
recovered by an uncharged random selector or a terminal decoder.
The central question of this paper is how to read that cost from the
saddle itself.

The underlying experiment consists of a preparation, positive normalized
kernels, a causal action interface, executable future tests, and charged
register cuts. The unknown parameter is fixed throughout an execution.
The preparation determines path laws and posterior continuations; no
geometric reference measure is supplied independently of the experiment.
The predictive quotient identifies histories that admit identical legal
future tests. Its stochastic realizations, rather than ordinary linear
rank alone, govern finite memory.

\subsection{Saddle geometry and its realizations}
For a finite causal interface, all controller wirings within a given
profile are represented by common normalized positive matrix products.
The generators at successive cuts must be compatible under every
one-step continuation. Independently factoring each cut does not suffice.
Let $\mathcal C_{\mathbf K}$ be this compact, generally nonconvex
response set, and let $U$ be the unrestricted minimax value. A
least-favorable prior exposes the Bayes face $F_*$. The profiles attaining
$U$ are precisely those admitting compatible positive realizations of
\[
 \mathcal S_*=F_*\cap\{C:g_\theta(C)\ge U\text{ for all }\theta\}.
\]
Theorem~\ref{thm:saddle-realization} gives the exact loss decomposition,
attainment criterion, supporting-parameter contact equations, and the
factorizations imposed at a retained-state cut. This is not a minimax
interchange over nonconvex machines. Its operative distinction is between
a Bayes exposed face and its uniformly optimal part.

This distinction is decisive when some Bayes coefficients vanish.
Contact and stationarity at interior supporting parameters may determine
these coordinates even though the Bayes objective does not. Once they
are determined, a convex-face obstruction applies to every stochastic
encoder and decoder. Lemma~\ref{lem:face-packing} gives a simple reusable
form: obligatory vertices and additional rows in disjoint faces require
separate positive generators. Exact causal recodings preserve the
attainable profiles; approximate simulations transport the value with
explicit state products and error allowances.

\subsection{Exact memory in a continuous marked experiment}
Our main application is
\[
 b(p,q)=\Ber(p)\otimes\Ber(q)\otimes\Ber(1/2),\qquad(p,q)\in[0,1]^2,
\]
where the last coordinate is the actual preparation mark. The target
$Q_\gamma$ has masses
$((1+\gamma)/4,(1-\gamma)/4,(1-\gamma)/4,(1+\gamma)/4)$ on
$000,001,110,111$, and zero mass elsewhere. After candidate training,
the auditor freezes a row and event and applies them to independent
candidate and target validations. The score is target indicator minus
candidate indicator. Write $U_N(\gamma)$ for its unrestricted minimax
value with $N$ training preparations.

The exact two-preparation saddle of Theorem~\ref{thm:two-family} has a
two-point least-favorable prior for $6/25\le\gamma\le13/50$.
Its two extreme-count response coordinates are Bayes ties.
Theorem~\ref{thm:five-memory} proves that every minimax response, not
only the displayed symmetric optimizer, must average these coordinates
to the same explicitly determined number $t_\gamma\in(0,1)$.
Three other count rows are forced cube vertices. Two disjoint exposed
faces then require two additional generators. Consequently
\[
 \boxed{\text{minimum decision width at the exact two-preparation value}=5.}
\]
The result applies to the original adaptive row--event interface:
training, gates, stopping, retained randomness, and independent calibration
are included. A common-garbling subfamily supplies the converse; the
construction works against the original private candidate class.

Applying the same argument to future \emph{output channels} at a
validation cut, rather than only to event responses, yields a stronger
scheduled result. If the candidate's complete validation is read before
any target validation report, eight deterministic continuation vertices
and four additional disjoint faces force twelve generators. Theorem~\ref{thm:serial-twelve}
therefore proves
\[
 \boxed{\text{minimum peak for exact forward serial validation}=12.}
\]
The matching implementation is the existing five-event residual machine.
The lower bound does not assume that the event name remains stored after
its residual is sufficient. This serial-order assertion is distinguished
from the order-independent decision-width theorem. Reverse-order and
interleaved validations are not assigned a peak optimum by this result.

For the collision bias $\gamma=1/4$, the exact statistical values are
\begin{equation}\label{eq:intro-exact-values}
 U_1=\frac{85-7\sqrt{73}}{64},\qquad
 U_2=\frac{-r^3+6r^2-3r+14}{32},\quad 2r^3+25r^2-3=0,
\end{equation}
where $1/3<r<7/20$. Their full-parameter lower certificates and matching
Bayes upper bounds are proved in the article. The new memory converses
extract consequences from the entire optimal face of these calculations.
They do not replace the exact values by a relaxation or a finite grid.
Compactness also gives a strict four-state loss uniformly over the
bias interval, and a strictly positive eleven-state loss for the forward
serial class. The four-state loss survives sufficiently small positive
sensor noise through an explicit vanishing Gaussian-tail modulus.

\subsection{A law for every preparation number}
The finite-$N$ saddles also determine a natural limiting question.
Theorem~\ref{thm:prior-localization} proves that unrestricted values
increase to the distance from the target to the nearest candidate
experiment. More significantly, every least-favorable prior concentrates
on that closest-experiment set. This follows from the same common
empirical response used as a competitor in the saddle game; the prior
is not replaced by a parameter-dependent algorithm.

For the marked family the closest candidates and a uniform quadratic
growth inequality are derived directly from its four diagonal masses.
With $s_\gamma=\sqrt{(1+\gamma)/2}$, the exact limiting value is
\[
 d_\gamma=2s_\gamma-1-\gamma/2,
\]
and its contact set consists exactly of $(s_\gamma,s_\gamma)$ and
$(1-s_\gamma,1-s_\gamma)$. For every $N\ge1$,
\[
 d_\gamma-\sqrt{3/N}\le U_N(\gamma)\le d_\gamma.
\]
Every symmetric least-favorable prior converges to the equal mixture of
those points, with quadratic-transport distance at most
$\sqrt{2300/7}(3/N)^{1/4}$. The intrinsic growth constant is derived,
not assumed as a dimension or small-ball hypothesis. These are uniform
bounds and a limiting classification, not a claim that a two-point
prior is exact at every finite preparation number or that the rates
are sharp.

\subsection{Prepared foundations and retained consequences}
The general foundations layer proves path construction, predictive
quotient closure, checkpoint resolution, and acquired causal transport.
The latter jointly charges calibration, preparation, quantization,
numerical and model error, and the memory of a downstream controller.
The informative routing family exhibits the exact values
$1/(2\lceil m/K\rceil)$ with a real routing cut and $K/(2m)$ with a
free persistent controller selector. The separate nonlinear occupation
recursion retains all common-row constraints. These results distinguish
an unrestricted saddle from its legal positive realizations.

The positive-noise collision construction realizes the marked experiment
on the original report and actual-mark interface. The exact ideal score
for $N=2$, $W\ge12$, the finite-precision implementation, and the
transported confidence theorem remain part of the paper. So do the
integer revelation laws, compatible subdivision certificates, and the
explicitly attributed autonomous finite-memory comparison. For reference,
the deterministic precursor constants are
\begin{equation}\label{eq:constants}
 u_1=\frac{4035777}{10240000},\qquad
 m_2=\frac{12-3\sqrt3}{16},\qquad\beta_0=\frac1{300}.
\end{equation}
Their earlier constructions and proofs are retained along with their
later exact improvements.

\subsection{Relation to earlier work}
Positive stochastic realization and invariant cone methods have classical
antecedents in Heller \cite{Heller}; Vidyasagar \cite{Vidyasagar} surveys
the complete hidden-Markov realization problem. Nonnegative rank differs
from ordinary linear rank, and its geometric interpretation is developed
by Gillis and Glineur \cite{GG}. Our matrix normalization adds the
specified causal action order and charged cuts. We do not claim the
positive-realization mechanism itself as new. The additional argument
links its generators to the uniformly optimal portion of a statistical
saddle, including contact equations on Bayes-indifferent coordinates.

Blackwell comparison and convex minimax provide the statistical
mechanisms \cite{Blackwell,Sion}. Belief-state dynamic programming is
classical \cite{SmallwoodSondik}. Controlled information acquisition and
adaptive sensing have an established history
\cite{Chernoff,NJ,NAV}; neither causal action selection nor an adaptive
advantage by itself is a priority claim here. Finite-memory testing and
randomization are likewise classical \cite{HC,HCrandom,FH,CFH}. The
autonomous stationary formula below is attributed to Hellman and Cover.
Recent work on memory-constrained adversarial testing \cite{MP} and
imperfect-recall optimization \cite{ZSV} concerns adjacent constraints
with different adversarial or objective conventions. Polynomial policy
optimization and positivity methods also have substantial antecedents
\cite{MM,Putinar}. We retain the explicit distinctions in the literature
crosswalk accompanying the source.

The empirical estimation bound used in the preparation limit is the
same elementary finite-alphabet estimate appearing in the retained
decision-spectrum theorem. The new use is localization of every
least-favorable prior, followed by an explicit intrinsic contact
calculation. Norberg's filtered-experiment comparison \cite{Norberg}
remains an adjacent general theory; no proof-level non-overlap claim is
made without a complete comparison of the original work.

\subsection{Organization}
The prepared foundations and finite interface are followed by statistical
and controlled duality. Section~\ref{sec:saddle-realization} gives the
central realization criterion. The exact marked saddles and physical
bridge supply its hypotheses, and Section~\ref{sec:exact-memory}
extracts the decision and serial-peak converses. Section~\ref{sec:preparation-localization}
gives the all-preparation law. The remaining resource and physical
consumers follow. Complete predecessor volumes retain the development
of the predictive, geometric, and operator interfaces. The independent
historical A2 chain and the B4/C2 aggregate obligations are not turned
into consequences of a finite-state theorem by a change of terminology.
